\begin{solapartat}
\punts{0,25} L'equació de la posició vertical és la component del vector posició de la massa respecte del centre del disc que podem escriure en funció de l'angle $\theta$ respecte de l'eix vertical $y$. L'angle augmenta linealment amb el temps i proporcionalment a la velocitat angular $\theta = \omega t + \varphi_0$. I, per tant: $y(t) = A\cos(\theta) = A\cos(\omega t + \varphi_0)$ o alternativament amb el sinus: $y(t) = A\sin(\theta) = A\sin(\omega t + \varphi_0)$

\punts{0,25} Per escriure l'equació necessitem l'amplitud, que és el radi $A = 0,19\ \mathrm{m}$, la velocitat angular $\omega = 6,41\ \mathrm{rad/s}$ i la fase inicial l'obtenim de l'instant inicial $t = 0$ en què la posició és $-A$

\punts{0,25} Càlcul de la posició:

Si es considera el sentit positiu cap amunt:

$y(0) = -A = A\cos(\varphi_0)$ \ i \ $-1 = \cos(\varphi_0)$, $\varphi = \arccos(-1) = \pi\ \mathrm{rad}$

I, per tant, $y(t) = 0,19\cos(6,41\,t + \pi)\ \mathrm{m}$ i $t$ en segons,

o també $y(t) = -0,19\cos(6,41\,t)\ \mathrm{m}$ i $t$ en segons.

Si es considera el sentit positiu cap avall:

$y(0) = A = A\cos(\varphi_0)$ \ i \ $1 = \cos(\varphi_0)$, $\varphi = \arccos(1) = 0\ \mathrm{rad}$

I, per tant, $y(t) = 0,19\cos(6,41\,t)\ \mathrm{m}$ i $t$ en segons,

o també $y(t) = -0,19\cos(6,41\,t + \pi)\ \mathrm{m}$ i $t$ en segons.

\punts{0,25} Les velocitats angulars de les dues masses són iguals; per tant, $\omega = \sqrt{\frac{k}{m}}$ i $k = \omega^2 m = 6,41^2\cdot 0,5 = 20,54\ \mathrm{N/m}$.

\punts{0,25} L'energia mecànica és $E_m = \frac{1}{2}kA^2 = 0,5\cdot 20,54\cdot 0,19^2 = 0,371\ \mathrm{J}$
\end{solapartat}

\begin{solapartat}
\punts{0,25} Càlcul de la velocitat:

Si es considera el sentit positiu cap amunt:
% A l'original hi diu "m/si" (probablement "m/s i").
\[ v(t) = \frac{dy}{dt} = -0,19\cdot 6,41\sin(6,41\,t + \pi) = -1,218\sin(6,41\,t + \pi)\ \mathrm{m/s} \]

Si es considera el sentit positiu cap avall:
\[ v(t) = \frac{dy}{dt} = -0,19\cdot 6,41\sin(6,41\,t) = -1,218\sin(6,41\,t)\ \mathrm{m/s} \]

\punts{0,25} L'energia cinètica:
\[ E_c = \tfrac{1}{2}mv^2 = \tfrac{1}{2}0,5\cdot 1,218^2\sin^2(6,41\,t + \pi) = 0,371\sin^2(6,41\,t + \pi)\ \mathrm{J} \]
o bé
\[ E_c = \tfrac{1}{2}mv^2 = \tfrac{1}{2}0,5\cdot 1,218^2\sin^2(6,41\,t) = 0,371\sin^2(6,41\,t)\ \mathrm{J} \]

Gràfica:

$E_m = 0,371\ \mathrm{J}$, \quad $E_p = \frac{1}{2}ky^2 = 10,27y^2\ \mathrm{J}$,

$E_c = E_m - E_p = 0,371 - 10,27y^2\ \mathrm{J}$

\figura[0.45]{sol_fig1.png}

\punts{0,25} Representació de l'$E_m$

\punts{0,25} Representació de l'$E_p$

\punts{0,25} Representació de l'$E_c$

Si l'escalatge no és correcte, s'han de restar 0,25 punts.
\end{solapartat}
